\section*{Chapter \ref{ch:b2:coordination-complexes} --- Transition Metals and Coordination Complexes}

\begin{solution}{exo:b2:coordination-complexes:1}
\ce{Ti^3+} $d^1$, \ce{V^3+} $d^2$, \ce{Cr^3+} $d^3$, \ce{Mn^2+} $d^5$, \ce{Fe^3+} $d^5$, \ce{Co^2+}
$d^7$, \ce{Ni^2+} $d^8$, \ce{Zn^2+} $d^{10}$.
\end{solution}

\begin{solution}{exo:b2:coordination-complexes:2}
Water 1, en 2, ethanedioate 2, edta 6, thiocyanate 1; thiocyanate is ambidentate
(through S or N).
\end{solution}

\begin{solution}{exo:b2:coordination-complexes:3}
Tetraamminecopper(II) sulfate; potassium hexacyanidoferrate(III);
tetraamminedichloridocobalt(III) chloride; tetracarbonylnickel(0).
\end{solution}

\begin{solution}{exo:b2:coordination-complexes:4}
Linear; tetrahedral; square planar ($d^8$, third series); octahedral.
\end{solution}

\begin{solution}{exo:b2:coordination-complexes:5}
Two: cis (the two chlorides adjacent) and trans (opposite). In a tetrahedron there
would be only one.
\end{solution}

\begin{solution}{exo:b2:coordination-complexes:6}
\ce{Cr(CO)6} 18; \ce{Fe(CO)5} 18; \ce{Ni(CO)4} 18; \ce{V(CO)6}: $5 + 12 = 17$;
\ce{Mn2(CO)10} 18 per manganese. \ce{V(CO)6} breaks the rule: it is \omterm{def:b2:diatomic-mos:magnetism}{paramagnetic} and
easily reduced to \ce{[V(CO)6]-}, which has 18.
\end{solution}

\begin{solution}{exo:b2:coordination-complexes:7}
Ferrocene: Fe(II) $d^6$ + 12 = 18. Cobaltocene: Co(II) $d^7$ + 12 = 19. Cobaltocene loses
its extra electron easily, giving the cobaltocenium ion \ce{[Co(C5H5)2]+}, with 18.
\end{solution}

\begin{solution}{exo:b2:coordination-complexes:8}
One edta replaces six water molecules: \ce{[Ni(H2O)6]^2+ + edta^4- -> [Ni(edta)]^2- +
6H2O}, seven particles from two. The entropy of reaction is strongly positive, and
$K$ very large. Six separate ligands replace six water molecules with no gain in the
number of particles.
\end{solution}

\begin{solution}{exo:b2:coordination-complexes:9}
Nitrito-$\kappa N$ (\ce{Co-NO2}, the nitro isomer, yellow) and nitrito-$\kappa O$
(\ce{Co-O-N=O}, the nitrito isomer, red).
\end{solution}

\begin{solution}{exo:b2:coordination-complexes:10}
Three: the trans isomer (achiral: mirror planes) and the two enantiomers of the cis
isomer (the two chelates and the two chlorides wind one way or the other).
\end{solution}

\begin{solution}{exo:b2:coordination-complexes:11}
Planar hexagon: B at positions 1,2 (ortho), 1,3 (meta), 1,4 (para): three isomers.
Trigonal prism: B along an edge of a triangle, along a vertical edge, or across a
rectangular face: three. The octahedron gives two. Finding no more than two, despite
many attempts and several compounds, rules out the other two geometries --- if no
third isomer has been overlooked.
\end{solution}

\begin{solution}{exo:b2:coordination-complexes:12}
\ce{[RhCl(PPh3)3]}: Rh(I) $d^8$ + 4 × 2 = 16. \ce{[IrH(CO)(PPh3)3]}: Ir(I) $d^8$ + 5 × 2 =
18. The rhodium complex, with 16 electrons, can take two more (oxidative addition,
\cref{ch:b2:catalytic-cycles}).
\end{solution}

\begin{solution}{pb:b2:coordination-complexes:1}
\textbf{1.} The chlorides outside the \omterm{def:b2:coordination-complexes:coordination-sphere}{coordination sphere}, free in solution.
\textbf{2.} A: 0; B: 1; C and D: 2.
\textbf{3.} A: $[\ce{Co(NH3)6}]^{3+}$ + 3 \ce{Cl-}: 4 ions; B: 3; C, D: 2. They agree.
\textbf{4.} +3.
\textbf{5.} $d^6$.
\textbf{6.} Six in each: 6 \ce{NH3}; 5 \ce{NH3} + 1 Cl; 4 \ce{NH3} + 2 Cl.
\textbf{7.} \ce{[Co(NH3)6]Cl3}; \ce{[CoCl(NH3)5]Cl2}; \ce{[CoCl2(NH3)4]Cl} (twice).
\textbf{8.} Hexaamminecobalt(III) chloride; pentaamminechloridocobalt(III) chloride.
\textbf{9.} Tetraamminedichloridocobalt(III).
\textbf{10.} Geometric (cis--trans) isomers.
\textbf{11.} \ce{[CoCl3(NH3)3]}, a neutral complex: no ions, no chloride precipitated at
once.
\textbf{12.} Two: fac and mer.
\textbf{13.} Cis: the two Cl at $90^\circ$; trans: at $180^\circ$.
\textbf{14.} Three (ortho, meta, para positions on the hexagon).
\textbf{15.} Three.
\textbf{16.} It supports the octahedron, the only one of the three that predicts two.
\textbf{17.} No: absence of a third isomer is negative evidence. A third isomer would
have ruled out the octahedron.
\textbf{18.} The one that forms the chelate: a bidentate ethanedioate can span only two
cis positions.
\textbf{19.} Three en chelates around cobalt, each spanning two cis positions.
\textbf{20.} Neither: no mirror plane and no centre of inversion.
\textbf{21.} Two enantiomers, $\Delta$ and $\Lambda$.
\textbf{22.} That the six ligands are arranged in three dimensions, as in an
octahedron; a planar arrangement would give an achiral ion.
\textbf{23.} Cis: chiral (two enantiomers). Trans: achiral.
\textbf{24.} It showed that optical activity belongs to molecular geometry in general,
not to carbon, and confirmed the octahedron.
\textbf{25.} The octahedron predicts $\boldsymbol{2}$ isomers of \ce{[CoCl2(NH3)4]+}; the
planar hexagon and the trigonal prism predict $\boldsymbol{3}$.
\end{solution}
